\documentclass{resume}

\usepackage[left=0.4in,top=0.4in,right=0.4in,bottom=0.4in]{geometry}
\usepackage{fontawesome5}
\newcommand{\tab}[1]{\hspace{.2667\textwidth}\rlap{#1}}
\newcommand{\itab}[1]{\hspace{0em}\rlap{#1}}
\name{Kevin Lee}
\address{+1 530-376-9732 \\ Davis, CA}
\address{\faEnvelope\ \href{mailto:yhylee@ucdavis.edu}{yhylee@ucdavis.edu} \quad
\faLinkedin\ \href{https://www.linkedin.com/in/kevin-lee-b78448388/}{linkedin.com/in/kevin-lee-b78448388} \quad
\faGithub\ \href{https://github.com/section9-us}{github.com/section9-us}}

\begin{document}
\small

\begin{rSection}{PROFESSIONAL SUMMARY}
Network-focused Software Engineer with 10+ years of experience assessing, securing, and supporting complex networked systems in public-sector and critical-infrastructure environments. Strong background in network vulnerability analysis, penetration testing, red teaming, packet analysis, incident response, and remediation consulting for nuclear, thermal power, railway control, and power-system networks. Combines network security expertise with software engineering and database experience to improve reliability and defensive visibility.
\end{rSection}

\begin{rSection}{EXPERIENCE}

\textbf{Software Engineer, Threat Management Division} \hfill Jan 2022 -- Present\\
Ministry of Science and ICT \hfill \textit{Sejong, Republic of Korea}
\begin{itemize}
    \itemsep -4pt {}
    \item Delivered 500+ CTI cases per year by analyzing malware, network attack evidence, indicators of compromise, and APT activity, supporting incident response and defensive remediation for government agencies.
\end{itemize}

\textbf{Software Engineer, Security Planning Division} \hfill May 2018 -- Jan 2022\\
Ministry of Science and ICT \hfill \textit{Sejong, Republic of Korea}
\begin{itemize}
    \itemsep -4pt {}
    \item Conducted network and application penetration testing, red teaming, and vulnerability analysis for 50+ public institutions annually (200+ total), including critical infrastructure networks supporting nuclear, thermal power, railway control, and power systems.
\end{itemize}

\textbf{Software Engineer, Reactor Design Development Division} \hfill Nov 2014 -- Mar 2018\\
KEPCO E\&C \hfill \textit{Gimcheon, Republic of Korea}
\begin{itemize}
    \itemsep -4pt {}
    \item Developed reactor design software and engineering databases for nuclear plant workflows while supporting vulnerability assessment and secure remediation for interconnected engineering systems.
\end{itemize}

\end{rSection}

\begin{rSection}{SKILLS}
\begin{tabular}{ @{} >{\bfseries}l @{\hspace{2ex}} p{0.77\textwidth} }
Programming & Python, Java, C/C++, C\#, JavaScript, HTML/CSS, SQL, Spring MVC \\
Networking & Network Security, Packet Analysis, Network Penetration Testing, Vulnerability Assessment, Red Teaming, Incident Response, Critical Infrastructure Networks, and Remediation \\
Security & Malware Analysis, CTI, Threat Modeling, Technical Security Assessments, and Secure Coding \\
Tools & Wireshark, Nmap, Burp Suite, IDA Pro, MySQL, PostgreSQL, Oracle, Docker, PyTorch, TensorFlow, Scikit-learn, and Hugging Face \\
Languages & Korean (Native), English (Fluent) \\
\end{tabular}\\
\end{rSection}

\begin{rSection}{EDUCATION}

{\bf Master of Science in Computer Science}, University of California, Davis (Current GPA: 4.0/4.0) \hfill {Expected Jun 2027}\\
Korean Government Long-term Fellowship for Overseas Study, Ministry of Science and ICT

{\bf Bachelor of Science in Computer Software}, Kwangwoon University (GPA: 3.32/4.0) \hfill {Feb 2015}

\end{rSection}

\begin{rSection}{CERTIFICATIONS \& TRAINING}
\textbf{Certifications:} Engineer Information Processing, HRDK (2013) \quad SQL Developer, Korea Data Agency (2014) \\
\textbf{Training:} Locked Shields, NATO CCDCOE (2023) \quad Advanced Web Application Exploits, Black Hat Training USA (2021)
\end{rSection}

\end{document}
